\chapter{Suites arithmétiques, géométriques et sommes}
\section*{Exercice 5.1 -- Suite arithmétique}
Une suite arithmétique satisfait $u_3=11$ et $u_8=26$. Trouver sa raison, $u_0$ et $\sum_{k=0}^8u_k$.

\section*{Exercice 5.2 -- Suite géométrique}
Pour $u_n=5(0{,}8)^n$, exprimer $\sum_{k=0}^{n}u_k$ et sa limite intuitive.

\section*{Exercice 5.3 -- Une raison négative}
Pour $u_n=(-2)^n$, calculer $u_0+\cdots+u_4$. La suite est-elle monotone ?

\section*{Exercice 5.4 -- Somme décalée}
Calculer $\sum_{k=4}^{12}(3k-1)$ en justifiant le nombre de termes.

\section*{Exercice 5.5 -- Télescopage}
Prouver $\sum_{k=1}^n k^2=n(n+1)(2n+1)/6$ en sommant $(k+1)^3-k^3$.

\section*{Exercice 5.6 -- Somme des cubes}
En posant $T_k=k(k+1)/2$, prouver $\sum_{k=1}^nk^3=T_n^2$.
